
3. \underline{Geometric proof of the monodromy theorem}

Let $\Lambda$ be as in \SGARefLink{sga7.sga7i.ei.n2-2}{2.2}.

\underline{Conjecture} \SGARefTarget{sga7.sga7i.ei.n3-1}{3.1} (\underline{purity}). \underline{Let $A$ be an
excellent regular strictly local ring} (\SGARefLink{sga7.sga7i.ei.n0-0-2}{0.0.2}) \underline{and let $D$ be a
regular divisor in $\Spec(A)$ defined by a parameter $x$. Put
$U=\Spec(A)-D=\Spec(A[1/x])$. Then}

\begin{equation}\tag{\SGARefTarget{sga7.sga7i.ei.n3-1-1}{3.1.1}}
H^{q}(U,\Lambda)=
\begin{cases}
\Lambda & \text{if }q=0,\\[4pt]
\Lambda(-1) & \text{if }q=1\quad\text{(notation (\SGARefLink{sga7.sga7i.ei.n0-0-6}{0.0.6})),}\\[4pt]
0 & \text{if }q>1.
\end{cases}
\end{equation}

Here are some cases in which this conjecture is known (SGA 4 XIX): for
$q=0$ or $1$; when $A$ has characteristic $0$; when $A$ has equal
characteristic $p$ and resolution of singularities is available in
characteristic $p$ in dimensions $\leqslant\dim(A)$; or when
$(\Spec(A),D)$ is a smooth pair (the relative purity theorem, SGA 4 XVI
3.7).

When purity is available and
$D=\sum_{i\in B}D_i$ is a normal-crossings divisor in $\Spec(A)$ defined
by part of a regular system of parameters, the cohomology of
$U=\Spec(A)-D$ is given by


\begin{equation}\tag{\SGARefTarget{sga7.sga7i.ei.n3-1-2}{3.1.2}}
\begin{aligned}
H^0(U,\Lambda)&=\Lambda,\\[4pt]
H^1(U,\Lambda)&=\Lambda(-1)^B,\\[4pt]
H^q(U,\Lambda)&=\bigwedge\nolimits^q H^1(U,\Lambda)
  \qquad\text{(via the cup product).}
\end{aligned}
\end{equation}
These formulas are identical to those which, over $\mathbf C$, give the
cohomology of $D^{*B}\times D^k$.

\SGARefTarget{sga7.sga7i.ei.n3-2}{3.2}. Let $S$ be a strictly local trait, and retain the notation of
(\SGARefLink{sga7.sga7i.ei.n0-0-1}{0.0.1}) and (\SGARefLink{sga7.sga7i.ei.n0-0-3}{0.0.3}). Let $f:X\longrightarrow S$ be a morphism of finite
type. Suppose that $X$ is regular, that $X_\eta$ is smooth, and that
$X_s$ is a normal-crossings divisor in $X$. We shall suppose that this
divisor is, globally and not merely locally for the \'{e}tale topology, a
sum of regular divisors:
\[
(X_s)_{\mathrm{red}}=\sum_{i\in B}D_i
\]
(the general case can be recovered by localization). Denote their
multiplicities by $m_i$, so that
$X_s=\sum_{i\in B}m_iD_i$.

Let $x$ be a point of $X_s$, let $\bar x$ be a geometric point of $X$
lying over $x$ (for example, $x$ may be closed and equal to $\bar x$),
put $C=\{i\mid x\in D_i\}$, and let $K$ be the complex concentrated in
degrees $0$ and $-1$
\begin{equation*}
K:\mathbf Z^C\xrightarrow{\ d\ }\mathbf Z,
\qquad d((n_j))=\sum_jm_jn_j.
\end{equation*}

\underline{Theorem} \SGARefTarget{sga7.sga7i.ei.n3-3}{3.3} (\underline{uses purity}). \underline{The groups
of tame vanishing cycles $\psi^i_{t\bar x}$ are as follows:}

(i) \underline{There is a canonical isomorphism of graded
$\Lambda$-algebras}
\begin{equation*}
\psi^*_{t\bar x}\simeq
\bigwedge\nolimits^*\bigl(H_1(K)\otimes\Lambda(-1)\bigr)
\otimes_\Lambda\psi^0_{t\bar x}.
\end{equation*}

(ii) $\psi^0_{t\bar x}$ \underline{is the} $\Lambda$-\underline{algebra}
$\Lambda^E$, \underline{where $E$ is a set on which the tame inertia
group} $I_t=I/P=\widehat{\mathbf Z}(1)(k(s))$
\underline{acts transitively, and whose cardinality is the largest
prime-to-$p$ divisor of $\#H_0(K)$.}

Here is the method of proof.

a) Let $X_{(\bar x)}$ be the strict localization of $X$ at $\bar x$,
let $U=X_{(\bar x)}-X_{(\bar x)s}$ (the complement of a
normal-crossings divisor), let $t_j$ be a local equation for $D_j$, put
\[
X_n=X_{(\bar x)}\bigl[(t_j^{1/n})\bigr]
\qquad ((n,p)=1),
\]
which is a regular scheme, and let $U_n$ be the inverse image of $U$ in
$X_n$. By (\SGARefLink{sga7.sga7i.ei.n3-1-2}{3.1.2}),
\[
H^q(U_n,\Lambda)=\bigwedge\nolimits^q\bigl(\Lambda(-1)^C\bigr).
\]
For $m=nd$, the pullback map from $H^q(U_n,\Lambda)$ to
$H^q(U_m,\Lambda)$ is multiplication


by $d^q$. Putting
$\widetilde U=\varprojlim_{(n,p)=1}U_n$, one therefore has
\begin{equation}\tag{\SGARefTarget{sga7.sga7i.ei.n3-3-1}{3.3.1}}
\begin{cases}
H^0(\widetilde U,\Lambda)=\Lambda,\\[6pt]
H^q(\widetilde U,\Lambda)=\displaystyle\varinjlim_n
  H^q(U_n,\Lambda)=0 & \text{if }q\ne0.
\end{cases}
\end{equation}

b) $\widetilde U$ is a pro-covering of $U$ with group
$\widehat{\mathbf Z}(1)(k(s))^C$. By definition,
\[
\psi^*_{t\bar x}=H^*(U_{\eta_t},\Lambda).
\]
$\widetilde U$ is a pro-covering, with group
$H_1(K)\otimes\widehat{\mathbf Z}(1)(k(s))$, of each connected component
of $U_{\eta_t}$, and these components are permuted transitively by the
tame inertia group. Formula (\SGARefLink{sga7.sga7i.ei.n3-3-1}{3.3.1}) and the Hochschild--Serre spectral
sequence now give the formulas in \SGARefLink{sga7.sga7i.ei.n3-3}{3.3}.

Pictorially, one may say that $U_{\bar\eta}$ is the fibre of a morphism
$\delta$ from the cohomological
$K(\widehat{\mathbf Z}(1)(k(s))^C,1)$ represented by $U$ to the
cohomological $K(\widehat{\mathbf Z}(1)(k(s)),1)$ represented by $\eta$;
the morphism $\delta$ induces
\[
d:(n_j)\longmapsto\sum_jm_jn_j:
\widehat{\mathbf Z}(1)(k(s))^C
\longrightarrow\widehat{\mathbf Z}(1)(k(s)).
\]

\underline{Corollary} \SGARefTarget{sga7.sga7i.ei.n3-4}{3.4} (\underline{uses purity}). \underline{Suppose
that $f$ is proper. Let $N$ be the least common multiple of the
multiplicities $m_j$ $(j\in B)$. Let $q\geqslant0$, and let $q'$ be the
minimum of $q$ and the greatest number of divisors $D_j$ passing through
one point. For $T$ in the tame inertia group $I_t=I/P$, one has}
\begin{equation*}
(T^N-1)^{q'+1}=0
\quad\underline{\text{on}}\quad
H^q(X_{\bar\eta},\Lambda)^P.
\end{equation*}
This follows at once from \SGARefLink{sga7.sga7i.ei.n3-3}{3.3}, Scholium \SGARefLink{sga7.sga7i.ei.n2-5}{2.5}, and Variant \SGARefLink{sga7.sga7i.ei.n2-7}{2.7}.

The required purity results are available for $H^1$, in characteristic
$0$, or when $S$ is essentially of finite type over a field.

\underline{Theorem} \SGARefTarget{sga7.sga7i.ei.n3-5}{3.5}. \underline{Let $C$ be a smooth projective curve
over the fraction field $k(\eta)$ of a discrete valuation ring $V$. The
inertia group $I$ acts on $H^1(C_{\bar\eta},\mathbf Q_\ell)$ by
quasi-unipotent automorphisms of level $2$: there exists $N$ such that,
for $T\in I$,}
\begin{equation*}
(T^N-1)^2=0
\quad\text{on}\quad H^1(C_{\bar\eta},\mathbf Q_\ell).
\end{equation*}
We reduce to the case in which $V$ is strictly local. The assertion is
that there exists an open subgroup $I_1$ of $I$ such that
$(T-1)^2=0$ on $H^1$ for $T\in I_1$.


We may extend the scalars of $k(\eta)$ to a finite extension, replacing
$V$ by its normalization in that extension. Since the image of the
pro-$p$ group $P$ in an $\ell$-adic Lie group
$\mathrm{GL}(n,\mathbf Q_\ell)$ is finite, we may in particular suppose
that $P$ acts trivially.

Let $C_o$ be a minimal model of $C$ over $V$ (for the existence of
$C_o$, which is based on Abhyankar's theorem, see the discussion in [3],
\S~2, p.~87; one could instead reduce to the case where $V$ is excellent
by passing to its completion and invoking SGA 4 XVI 1.6).

Let $C_1$ be a regular model of $C$, obtained from $C_o$ by blowups, whose
special fibre is a normal-crossings divisor. The arguments of \SGARefLink{sga7.sga7i.ei.n3-3}{3.3} and
\SGARefLink{sga7.sga7i.ei.n3-4}{3.4} apply to $C_1$, $\psi_t^0$, $\psi_t^1$, and
\[
H^1(C_{\bar\eta},\mathbb Z/\ell^n)^P
=H^1(C_{\bar\eta},\mathbb Z/\ell^n),
\]
with $N$ independent of $n$; Theorem \SGARefLink{sga7.sga7i.ei.n3-5}{3.5} follows.

\underline{Remark} \SGARefTarget{sga7.sga7i.ei.n3-6}{3.6}. Let $A$ be an abelian variety over $k(\eta)$.
Since $A$ is a quotient of a Jacobian, it follows again that the action
of $I$ on $H^1(A_{\bar\eta},\mathbb Q_\ell)$, or on $T_\ell(A)$, is
quasi-unipotent of level $2$.

\underline{Remark} \SGARefTarget{sga7.sga7i.ei.n3-7}{3.7}. The action of $I$ on
$H^1(X_{\bar\eta},\mathbb Q_\ell)$ is in fact quasi-unipotent of level
$2$ for every scheme $X$ of finite type over $\eta$.

Theorem \SGARefLink{sga7.sga7i.ei.n3-5}{3.5} (or Remark \SGARefLink{sga7.sga7i.ei.n3-6}{3.6}) is the basis for the proof of the stable
reduction theorem for abelian varieties that will be given in Expos\'{e}
IX, \S~3.

\SGARefTarget{sga7.sga7i.ei.n3-8}{3.8}. It is clear that whenever purity and resolution of singularities
are available (for example, in characteristic $0$), the preceding method
proves the monodromy theorem and gives bounds for the nilpotence
exponent. Thus, in equal characteristic $0$, if $X$ is proper and smooth
over $\eta$ and $T$ belongs to an open subgroup $I_1$ of $I$, one has
\begin{equation*}
(T-1)^{q+1}=0
\quad\text{on}\quad H^q(X_{\bar\eta},\mathbb Z_\ell)
\qquad(T\in I_1\subset I).
\end{equation*}
